
\documentclass{article}
\usepackage{colortbl}
\usepackage{makecell}
\usepackage{multirow}
\usepackage{supertabular}

\begin{document}

\newcounter{utterance}

\centering \large Interaction Transcript for game `imagegame', experiment `random\_grids', episode 13 with Qwen3{-}32B{-}t0.0.
\vspace{24pt}

{ \footnotesize  \setcounter{utterance}{1}
\setlength{\tabcolsep}{0pt}
\begin{supertabular}{c@{$\;$}|p{.15\linewidth}@{}p{.15\linewidth}p{.15\linewidth}p{.15\linewidth}p{.15\linewidth}p{.15\linewidth}}
    \# & $\;$A & \multicolumn{4}{c}{Game Master} & $\;\:$B\\
    \hline

    \theutterance \stepcounter{utterance}  
    & & \multicolumn{4}{p{0.6\linewidth}}{
        \cellcolor[rgb]{0.9,0.9,0.9}{
            \makecell[{{p{\linewidth}}}]{
                \texttt{\tiny{[P1$\langle$GM]}}
                \texttt{Let us do the following task. The goal is to fill an empty grid that looks like this:} \\
\\ 
\texttt{▢ ▢ ▢ ▢ ▢} \\
\texttt{▢ ▢ ▢ ▢ ▢} \\
\texttt{▢ ▢ ▢ ▢ ▢} \\
\texttt{▢ ▢ ▢ ▢ ▢} \\
\texttt{▢ ▢ ▢ ▢ ▢} \\
\\ 
\texttt{A filled grid below is 5 by 5 and can look like this:} \\
\\ 
\texttt{▢ ▢ ▢ ▢ ▢} \\
\texttt{▢ ▢ E ▢ ▢} \\
\texttt{▢ ▢ ▢ ▢ ▢} \\
\texttt{▢ ▢ ▢ ▢ ▢} \\
\texttt{X X X X X} \\
\\ 
\texttt{I want you to describe this grid to me, step by step. You don't need to describe the empty squares, which are denoted with "▢". Only describe the location of letters in the grid. Then you wait for me to say "What is your next command?", and then you continue with the next step. Take the size of the grid into consideration while giving commands. You can give only one command at each turn but it can refer to single or multiple cells, rows, columns or even the whole grid at once. When you have described everything, you say "Command: DONE". Only give the command and omit everything else.} \\
\\ 
\texttt{For the filled grid above, here are the example steps.} \\
\\ 
\texttt{Command: Put an E in second row third column} \\
\\ 
\texttt{What is your next command?} \\
\texttt{Command: Fill the last row with X} \\
\\ 
\texttt{What is your next command?} \\
\texttt{Command: DONE} \\
\\ 
\texttt{Ok. Please do this for the following example, which is a 5 by 5 grid.} \\
\texttt{▢ ▢ ▢ D D} \\
\texttt{▢ ▢ ▢ ▢ D} \\
\texttt{▢ ▢ ▢ D ▢} \\
\texttt{▢ ▢ D D D} \\
\texttt{▢ ▢ ▢ ▢ ▢} \\
\texttt{What is your next command?} \\
            }
        }
    }
    & \\ \\

    \theutterance \stepcounter{utterance}  
    & \multicolumn{4}{p{0.6\linewidth}}{
        \cellcolor[rgb]{0.9,0.9,0.9}{
            \makecell[{{p{\linewidth}}}]{
                \texttt{\tiny{[P1$\rangle$GM]}}
                \texttt{<think>} \\
\texttt{Okay, let's see. The user provided a 5x5 grid with some letters filled in. My task is to describe the positions of the letters step by step, one command at a time. The example they gave used commands like "Put an E in second row third column" and "Fill the last row with X". So I need to follow that format.} \\
\\ 
\texttt{First, I need to look at the grid they provided:} \\
\\ 
\texttt{Row 1: ▢ ▢ ▢ D D} \\
\texttt{Row 2: ▢ ▢ ▢ ▢ D} \\
\texttt{Row 3: ▢ ▢ ▢ D ▢} \\
\texttt{Row 4: ▢ ▢ D D D} \\
\texttt{Row 5: ▢ ▢ ▢ ▢ ▢} \\
\\ 
\texttt{I need to identify all the letters and their positions. Let's go row by row.} \\
\\ 
\texttt{Row 1 has D in columns 4 and 5. So that's (1,4) and (1,5).} \\
\\ 
\texttt{Row 2 has D in column 5. So (2,5).} \\
\\ 
\texttt{Row 3 has D in column 4. (3,4).} \\
\\ 
\texttt{Row 4 has D in columns 3, 4, 5. (4,3), (4,4), (4,5).} \\
\\ 
\texttt{Row 5 is all empty.} \\
\\ 
\texttt{Now, the challenge is to give commands that place these D's. The example used specific positions and then a fill command. Since the D's} \\
            }
        }
    }
    & & \\ \\

    \theutterance \stepcounter{utterance}  
    & & & \multicolumn{2}{p{0.3\linewidth}}{
        \cellcolor[rgb]{0.9,0.9,0.9}{
            \makecell[{{p{\linewidth}}}]{
                \texttt{\tiny{[GM$|$GM]}}
                \texttt{Invalid instruction format} \\
            }
        }
    }
    & & \\ \\

\end{supertabular}
}

\end{document}
